% GENERATED FILE. Edit canonical lessons, then run npm run generate:books.
% canonical-source-hash: fnv1a64:fe32fae649a15b8e
% canonical-lessons: TE-C545-abhyasincu, TE-C545-kasarattu, TE-C545-guripettu, TE-C545-toyu, TE-C545-pandem

\chapter{Training}
\label{ch:te-training}
\hlchaptermodality{545}

\begin{chapteropening}
\textbf{By the end of this chapter:} \emph{I can talk about practice, workouts, aiming, shoving and bets.}

Complete the last lesson of chapter 545: I can talk about practice, workouts, aiming, shoving and bets.
\end{chapteropening}

\section[\texorpdfstring{\te{అభ్యసించు} (abhyasiñcu)}{అభ్యసించు (abhyasiñcu)}]{\te{అభ్యసించు} (abhyasiñcu) — to study closely; to practise}

\label{lesson:TE-C545-abhyasincu}



\begin{quote}
Before the new one: say the Telugu for to prevail, then the Telugu for running.
\end{quote}

\subsection*{You'll want to know: \te{అభ్యసించు}}
\textbf{\te{అభ్యసించు}} — \emph{abhyasiñcu} — \textquotedblleft{}to study closely; to practise\textquotedblright{}.

\begin{quote}
\te{అతను} \te{రోజూ} \te{రెండు} \te{గంటలు} \te{సంగీతం} \te{అభ్యసిస్తాడు}.
\end{quote}

\begin{quote}
\emph{atanu rōjū reṇḍu gaṇṭalu saṅgītaṁ abhyasistāḍu.}
\end{quote}

\begin{quote}
He practises music for two hours every day.
\end{quote}

Past (I …): \te{అభ్యసించాను}, \emph{abhyasiñcānu}.

Formal; \te{అభ్యాసం} is practice or exercise.

\begin{cousinweb}
\textbf{\te{అభ్యసించు}} (\emph{abhyasiñcu}) is built on Sanskrit \emph{abhyasana}, \textquotedblleft{}practice, exercise\textquotedblright{} (Monier-Williams), with the Telugu verb ending -\te{ఇంచు}.
\end{cousinweb}

\begin{grammarlens}[title={What you've built}]
One more word for describing what you lived and felt.
\end{grammarlens}

\subsection*{Your turn}
\begin{itemize}
\raggedright
  \item \emph{Say it:} \emph{abhyasiñcu}
  \item \emph{Say it:} \emph{abhyasiñcu}, once more
  \item \emph{Say it:} \emph{atanu rōjū reṇḍu gaṇṭalu saṅgītaṁ abhyasistāḍu.} — He practises music for two hours every day.
  \item \emph{Recall:} say the Telugu for to prevail, then the Telugu for running, then say \emph{abhyasiñcu} again
\end{itemize}

\begin{tcolorbox}[breakable,colback=teal!4,colframe=teal!35!black,title={Before you move on}]
What does \te{అభ్యసించు} mean? (\textquotedblleft{}To study closely; to practise\textquotedblright{}.) Say it once more.
\end{tcolorbox}
\section[\texorpdfstring{\te{కసరత్తు} (kasarattu)}{కసరత్తు (kasarattu)}]{\te{కసరత్తు} (kasarattu) — a workout; drill}

\label{lesson:TE-C545-kasarattu}



\begin{quote}
Before the new one: say the Telugu for running, then the Telugu for to study closely.
\end{quote}

\subsection*{You'll want to know: \te{కసరత్తు}}
\textbf{\te{కసరత్తు}} — \emph{kasarattu} — \textquotedblleft{}a workout; drill\textquotedblright{}.

\begin{quote}
\te{క్రీడాకారులు} \te{ఉదయమే} \te{కసరత్తు} \te{చేశారు}.
\end{quote}

\begin{quote}
\emph{krīḍākārulu udayamē kasarattu cēśāru.}
\end{quote}

\begin{quote}
The athletes worked out early in the morning.
\end{quote}

\te{కసరత్తు} \te{చేయు}, to work out.

\begin{cousinweb}
\textbf{\te{కసరత్తు}} (\emph{kasarattu}) came into Telugu through Hindustani \emph{kasrat}, \textquotedblleft{}exercise\textquotedblright{}, which is usually traced to Arabic.
\end{cousinweb}

\begin{grammarlens}[title={What you've built}]
One more word for describing what you lived and felt.
\end{grammarlens}

\subsection*{Your turn}
\begin{itemize}
\raggedright
  \item \emph{Say it:} \emph{kasarattu}
  \item \emph{Say it:} \emph{kasarattu}, once more
  \item \emph{Say it:} \emph{krīḍākārulu udayamē kasarattu cēśāru.} — The athletes worked out early in the morning.
  \item \emph{Recall:} say the Telugu for running, then the Telugu for to study closely, then say \emph{kasarattu} again
\end{itemize}

\begin{tcolorbox}[breakable,colback=teal!4,colframe=teal!35!black,title={Before you move on}]
What does \te{కసరత్తు} mean? (\textquotedblleft{}A workout; drill\textquotedblright{}.) Say it once more.
\end{tcolorbox}
\section[\texorpdfstring{\te{గురిపెట్టు} (guripeṭṭu)}{గురిపెట్టు (guripeṭṭu)}]{\te{గురిపెట్టు} (guripeṭṭu) — to aim}

\label{lesson:TE-C545-guripettu}



\begin{quote}
Before the new one: say the Telugu for to study closely, then the Telugu for a workout.
\end{quote}

\subsection*{You'll want to know: \te{గురిపెట్టు}}
\textbf{\te{గురిపెట్టు}} — \emph{guripeṭṭu} — \textquotedblleft{}to aim\textquotedblright{}.

\begin{quote}
\te{అతను} \te{బంతిని} \te{గోలు} \te{వైపు} \te{గురిపెట్టాడు}.
\end{quote}

\begin{quote}
\emph{atanu bantini gōlu vaipu guripeṭṭāḍu.}
\end{quote}

\begin{quote}
He aimed the ball at the goal.
\end{quote}

Past (I …): \te{గురిపెట్టాను}, \emph{guripeṭṭānu}.

\te{గురి}, aim, + \te{పెట్టు}, to put.

\begin{cousinweb}
\textbf{\te{గురిపెట్టు}} (\emph{guripeṭṭu}) is native Dravidian: DEDR 1847 lists Telugu \emph{guṟi}, \textquotedblleft{}an aim, an object aimed at\textquotedblright{}, with Tamil \emph{kuṟi}; \emph{peṭṭu}, \textquotedblleft{}to put\textquotedblright{}, makes the verb.
\end{cousinweb}

\begin{grammarlens}[title={What you've built}]
One more word for describing what you lived and felt.
\end{grammarlens}

\subsection*{Your turn}
\begin{itemize}
\raggedright
  \item \emph{Say it:} \emph{guripeṭṭu}
  \item \emph{Say it:} \emph{guripeṭṭu}, once more
  \item \emph{Say it:} \emph{atanu bantini gōlu vaipu guripeṭṭāḍu.} — He aimed the ball at the goal.
  \item \emph{Recall:} say the Telugu for to study closely, then the Telugu for a workout, then say \emph{guripeṭṭu} again
\end{itemize}

\begin{tcolorbox}[breakable,colback=teal!4,colframe=teal!35!black,title={Before you move on}]
What does \te{గురిపెట్టు} mean? (\textquotedblleft{}To aim\textquotedblright{}.) Say it once more.
\end{tcolorbox}
\section[\texorpdfstring{\te{తోయు} (tōyu)}{తోయు (tōyu)}]{\te{తోయు} (tōyu) — to shove; to push}

\label{lesson:TE-C545-toyu}



\begin{quote}
Before the new one: say the Telugu for a workout, then the Telugu for to aim.
\end{quote}

\subsection*{You'll want to know: \te{తోయు}}
\textbf{\te{తోయు}} — \emph{tōyu} — \textquotedblleft{}to shove; to push\textquotedblright{}.

\begin{quote}
\te{జనం} \te{గుంపులో} \te{ఎవరో} \te{నన్ను} \te{తోశారు}.
\end{quote}

\begin{quote}
\emph{janaṁ gumpulō evarō nannu tōśāru.}
\end{quote}

\begin{quote}
Someone in the crowd shoved me.
\end{quote}

Past (I …): \te{తోశాను}, \emph{tōśānu}.

Rougher than \te{నెట్టు}.

\begin{cousinweb}
\textbf{\te{తోయు}} (\emph{tōyu}) is native Dravidian: DEDR 3340 lists Telugu \emph{trōcu}, \textquotedblleft{}to push, to shove, to thrust\textquotedblright{}, the older form of this verb.
\end{cousinweb}

\begin{grammarlens}[title={What you've built}]
One more word for describing what you lived and felt.
\end{grammarlens}

\subsection*{Your turn}
\begin{itemize}
\raggedright
  \item \emph{Say it:} \emph{tōyu}
  \item \emph{Say it:} \emph{tōyu}, once more
  \item \emph{Say it:} \emph{janaṁ gumpulō evarō nannu tōśāru.} — Someone in the crowd shoved me.
  \item \emph{Recall:} say the Telugu for a workout, then the Telugu for to aim, then say \emph{tōyu} again
\end{itemize}

\begin{tcolorbox}[breakable,colback=teal!4,colframe=teal!35!black,title={Before you move on}]
What does \te{తోయు} mean? (\textquotedblleft{}To shove; to push\textquotedblright{}.) Say it once more.
\end{tcolorbox}
\section[\texorpdfstring{\te{పందెం} (pandeṁ)}{పందెం (pandeṁ)}]{\te{పందెం} (pandeṁ) — a bet; a race}

\label{lesson:TE-C545-pandem}



\begin{quote}
Before the new one: say the Telugu for to aim, then the Telugu for to shove.
\end{quote}

\subsection*{You'll want to know: \te{పందెం}}
\textbf{\te{పందెం}} — \emph{pandeṁ} — \textquotedblleft{}a bet; a race\textquotedblright{}.

\begin{quote}
\te{ఎవరు} \te{ముందు} \te{వస్తారో} \te{అని} \te{మేము} \te{పందెం} \te{వేసుకున్నాం}.
\end{quote}

\begin{quote}
\emph{evaru mundu vastārō ani mēmu pandeṁ vēsukunnāṁ.}
\end{quote}

\begin{quote}
We had a bet on who would arrive first.
\end{quote}

\te{పందెం} \te{వేయు}, to bet.

\begin{cousinweb}
\textbf{\te{పందెం}} (\emph{pandeṁ}) is native Dravidian: DEDR 3921 lists Telugu \emph{pandemu}, \textquotedblleft{}a wager; a race\textquotedblright{}, with Tamil \emph{pantayam}, \textquotedblleft{}a contest for a prize, a stake\textquotedblright{}.
\end{cousinweb}

\begin{grammarlens}[title={What you've built}]
One more word for describing what you lived and felt.
\end{grammarlens}

\subsection*{Your turn}
\begin{itemize}
\raggedright
  \item \emph{Say it:} \emph{pandeṁ}
  \item \emph{Say it:} \emph{pandeṁ}, once more
  \item \emph{Say it:} \emph{evaru mundu vastārō ani mēmu pandeṁ vēsukunnāṁ.} — We had a bet on who would arrive first.
  \item \emph{Recall:} say the Telugu for to aim, then the Telugu for to shove, then say \emph{pandeṁ} again
\end{itemize}

\begin{tcolorbox}[breakable,colback=teal!4,colframe=teal!35!black,title={Before you move on}]
What does \te{పందెం} mean? (\textquotedblleft{}A bet; a race\textquotedblright{}.) Say it once more.
\end{tcolorbox}
